%%%PAGE 271 rot=0 legibility=fair summary=上部贴有两道印刷题（短道速滑过弯支撑力；三块木板进入粗糙斜面段）及铅笔演算；下部为空腔引力场、动滑轮弹簧等效劲度系数、天体质量公式、太阳变黑洞密度、未平衡摩擦的1/a-M图像等笔记
%%%BLOCK id=271-01 ch=04 sec=水平面圆周·摩擦力提供向心力 kind=problem alt=02 cont=none y=0.00-0.19
% 贴纸印刷题（上缘被截断）；原稿上有铅笔划线：“力支撑冰面才能保证”处有划线，“过弯道滑行时的运动员”处有划线
\begin{problem}[贴纸印刷题]
\unclear{…}2022年\unclear{…}威、曲春雨、范可新、武大靖和张雨婷组成的中国队夺得北京冬奥会短道速滑男女2000米混合接力冠军，为中国体育代表团收获了北京冬奥会的首枚金牌。短道速滑运动员过水平弯道时常用手支撑冰面以防侧滑，某运动员质量为$75\unit{kg}$，某次过弯道时的半径为$25\unit{m}$，速率为$36\unit{km/h}$，冰刀与冰面间的动摩擦因数为$0.2$，手套与冰面间的动摩擦因数为$0.8$，重力加速度$g=10\unit{m/s^2}$。过弯道滑行时的运动员手脚距离相对半径可忽略，弯道滑行的过程视为一段圆周运动，则该运动员至少用多大的力支撑冰面才能保证不发生侧滑（\struck{B}C） % 括号内手写答案：“B”被斜线划去，旁写C

A.\ $150\unit{N}$\qquad B.\ $200\unit{N}$\qquad C.\ $250\unit{N}$\qquad D.\ $300\unit{N}$

手写铅笔批注（部分被印刷图片遮挡，数字难辨）：
\begin{itemize}
\item 第2行右端圈出：$10\unit{m/s}$
\item $0.2+0.8(\unclear{780}+N)=75\cdot\dfrac{\unclear{100}}{\unclear{30}}$，\ 其下一行写$0.8(\unclear{780}+N)$（括号部分划有下划线），再下一行$=300$ % CHECK: “780”疑为750；式子本身（0.2、0.8的组合）按原稿照录
\item 印刷图片上写有：$0.6$
\item 右缘：$7.5$（下有横线）；圈出的$300$；其下$\unclear{1000}$（被页边截断）
\end{itemize}
\end{problem}
%%%END
%%%BLOCK id=271-02 ch=06 sec=动能定理·变力做功（平均力法） kind=problem alt=03 cont=none y=0.19-0.44
% 贴纸印刷题（上缘被撕边截断）；原稿圈出“正确”及“三块长木板恰能完全进入粗糙的PA段”，“PA段粗糙程度处处相同”处有铅笔划线
\begin{problem}[贴纸印刷题]
\unclear{如}图所示，一固定的足够长斜面体倾斜角为$\theta$，其中PB段光滑，PA段粗糙程度处处相同，三块完全相同且长度为$L$的质量分布均匀的木板123靠在一起，从N点静止释放。已知$PN=L$，三块长木板恰能完全进入粗糙的PA段，重力加速度为$g$，每块的质量为$m$，下列说法正确的是（\unclear{BD}） % 括号后手写答案，难辨，疑为BD

\begin{itemize}
\item[A.] PA段动摩擦因数$\mu=\dfrac{4}{3}\tan\theta$ % 原稿：式子被斜线划去
\item[B.] \unclear{三}块长木板运动过程中的最大速度为$v_m=\dfrac{5\sqrt{2gL\sin\theta}}{4}$ % 原稿：式下有大“✓”形勾画；该行开头被铅笔字遮盖
\item[C.] 第1块长木板恰完全进入粗糙的PA段时，第1、2块长木板之间的作用力为$F_{12}=2mg\sin\theta$ % 原稿：$2mg\sin\theta$被斜线划去，其后另有一笔“9”形
\item[D.] 第2块长木板恰完全进入粗糙的PA段时，速度大小为$v_2=\dfrac{\sqrt{22gL\sin\theta}}{3}$ % 原稿：句下有大“✓”形勾画
\end{itemize}

%FIG p271_a 0.495 0.235 0.99 0.41
\notefig[0.6]{figures/p271_a.png}

\begin{solution}
手写铅笔批注（分散在题面空白处，按位置照录）：
\begin{itemize}
\item 左侧页边：$W_{\text{合}}=0$
\item $3mg\cdot4L\sin\theta=$ \quad 图左侧续写：$\dfrac{\mu\cdot3mg\cos\theta}{2}\cdot3L$
\item 方框内：$\mu=\dfrac{8}{3}\tan\theta$
\item 图右上：$9mgL=6\mu mgL\cos\theta$ % 原稿“6”写在μ的左下方；等式左侧疑漏$\sin\theta$
\item 图右：\#\quad $\mu=\dfrac{3}{2\cos\theta}$
\item 选项C右下：$\dfrac{16}{9}mg\sin\theta$；其右圈出：$\unclear{\dfrac{2m}{18}}$ % 难辨
\item 选项C、D之间左侧另有难辨铅笔字迹：\unclear{T、3L}
\item 右下：平均力法\ \unclear{（…）}
\end{itemize}
\end{solution}
\end{problem}
%%%END
%%%BLOCK id=271-03 ch=05 sec=万有引力·空腔（割补法） kind=knowledge alt=none cont=none y=0.44-0.53
%FIG p271_b 0.06 0.45 0.175 0.54
\notefig[0.2]{figures/p271_b.png}
空腔中加速度恒定\ 指向地心
%%%END
%%%BLOCK id=271-04 ch=02 sec=弹力·等效劲度系数 kind=knowledge alt=06 cont=none y=0.47-0.61
%FIG p271_c 0.555 0.485 0.725 0.61
\notefig[0.25]{figures/p271_c.png}
求系统等效劲度系数。图中标注：$2x\downarrow$，$2kx$；$\downarrow4kx$；$x\downarrow$。
\[
K_{\text{等效}}=4K
\]
%%%END
%%%BLOCK id=271-05 ch=05 sec=天体质量·密度 kind=formula alt=none cont=none y=0.575-0.65
%FIG p271_d 0.06 0.575 0.165 0.65
\notefig[0.2]{figures/p271_d.png}
\[
M=\frac{\pi^2(2R+ct)^3}{2GT^2}
\]
图中标注$R$与$\dfrac{ct}{2}$。
%%%END
%%%BLOCK id=271-06 ch=05 sec=宇宙速度·黑洞 kind=derivation alt=none cont=none y=0.67-0.77
太阳变成黑洞\ $\therefore V_{\mathrm{II}}=C$.
\[
\sqrt2\sqrt{\frac{GM}{R}}=C\qquad R=\left(\frac{3M}{4\rho\pi}\right)^{\frac13}\qquad\rho=\frac{3C^6}{32\pi G^3M^2}
\] % 原稿：从$\sqrt{GM/R}$的分母R画弧线箭头指向$R=(\cdots)^{1/3}$
%%%END
%%%BLOCK id=271-07 ch=03 sec=验证牛顿第二定律·误差分析 kind=method exp=1 alt=none cont=none y=0.78-0.97
%FIG p271_e 0.07 0.785 0.238 0.875
\notefig[0.25]{figures/p271_e.png}
未平衡$f$：
\[
\frac{1}{a}=\frac{1}{mg-\mu Mg}\cdot M
\]
%FIG p271_f 0.24 0.878 0.44 0.955
\notefig[0.25]{figures/p271_f.png}
若$m$未远小于$M$：
%FIG p271_g 0.48 0.845 0.69 0.955
\notefig[0.25]{figures/p271_g.png}
%%%END
%%%PAGEEND 271
